% concatenated from arxiv-v3.tex
\documentclass[a4paper,11pt]{amsart}

\usepackage[left = 3.5 cm, right = 3.5 cm]{geometry}
\usepackage{amsthm,amssymb,amsmath,amsfonts,mathrsfs,amscd,amsbsy,dsfont,verbatim}
\usepackage{stmaryrd}
\usepackage{graphicx}
\usepackage[all,cmtip]{xy}
\usepackage{longtable}
%\usepackage[notref,notcite]{showkeys}
\usepackage{marginnote}
\usepackage{comment}
%\usepackage{enumitem}
%\pagestyle{empty}
\usepackage{babel}
\usepackage{latexsym}
\usepackage{mathtools}
\usepackage[utf8]{inputenc}
\usepackage{xcolor}
\usepackage{fancyhdr}
%\usepackage[nottoc]{tocbibind}
\usepackage{indentfirst}
\usepackage{paralist}
%\usepackage[a-1b]{pdfx}
\usepackage{hyperref}
%\usepackage{graphicx}
%\usepackage{dcpic, pictex}
%\usepackage{etoolbox}
\usepackage{tikz}
\usetikzlibrary{cd}
\usepackage{enumerate}
%\usepackage{makeidx}
\usepackage{multicol}
\usepackage{comment}
\usepackage{todonotes}
\usepackage[all,cmtip]{xy}
\usepackage[title]{appendix}
\usepackage{url}


\newcommand{\nirvana}[1]{{\color{red} \sf \(\star\star\star\) Nirvana: [#1]}}
\newcommand{\ja}[1]{{\color{magenta}\sf \(\star\star\star\) Jessica: [#1]}}

\newtheorem{theorem}{Theorem}[section]
\newtheorem{proposition}[theorem]{Proposition}
\newtheorem{lemma}[theorem]{Lemma}
\newtheorem{corollary}[theorem]{Corollary}
\newtheorem{conjecture}[theorem]{Conjecture}
\newtheorem{example}[theorem]{Example}
\newtheorem{question}[theorem]{Question}
\newtheorem*{theorem*}{Theorem}

\theoremstyle{definition}
\newtheorem{definition}{Definition}[section]
\newtheorem{remark}[theorem]{Remark}
\newtheorem{notation}[definition]{Notation}
\newtheorem*{claim}{Claim}
\newtheorem{note}{Note}[section]


\newcommand{\rest}{\mathbin\restriction}
\newcommand{\N}{\mathbb N}
\newcommand{\Q}{\mathbb Q}
\newcommand{\Z}{\mathbb Z}
\newcommand{\R}{\mathbb R}
\newcommand{\F}{\mathbb F}
\newcommand{\p}{\mathfrak{p}}
\newcommand{\f}{\mathfrak{f}}
\newcommand{\e}{\mathfrak{e}}
\newcommand{\Cc}{\mathcal C}
\newcommand{\Mat}{\text{Mat}}
\newcommand{\GL}{\text{GL}}
\newcommand{\Id}{\text{Id}}
\DeclareMathOperator{\im}{Im}
\DeclareMathOperator{\rk}{\text{rk}}
\DeclareMathOperator{\srk}{\text{srk}}
\DeclareMathOperator{\Gal}{Gal}
\DeclareMathOperator{\Aut}{Aut}
\DeclareMathOperator{\End}{End}
\DeclareMathOperator{\disc}{disc}
\DeclareMathOperator{\Frob}{Frob}
\newcommand{\sm}{{\scriptstyle\rm sm}}
\newcommand{\cO}{\mathcal{O}}
\newcommand{\cA}{\mathcal{A}}
\newcommand{\cB}{\mathcal{B}}
\newcommand{\cC}{\mathcal{C}}
\renewcommand{\bar}{\overline}
\DeclareMathOperator{\Tors}{Tors}

\newcommand{\mat}[4]{\left(\begin{array}{cc}
#1 &#2\\
#3 &#4
\end{array}\right)}
\newcommand{\lagrange}[2]{\left( \dfrac{#1}{#2}\right)}

\title{Torsion points on \(\GL_2\)-type abelian varieties}
\date{\today}

\author[J. Alessandrì]{Jessica Alessandrì}
\address{Max Planck Institute for Mathematics\\
Vivatsgasse 7\\
53111 Bonn\\
Germany}
\email{alessandri@mpim-bonn.mpg.de}
\urladdr{https://sites.google.com/view/jessicaalessandri}
\author[N. Coppola]{Nirvana Coppola}
\address{Dipartimento di Matematica, Università di Padova\\
Via Trieste, 63\\
35121 Padova\\
Italy}
\email{nirvana.coppola@unipd.it}
\urladdr{https://sites.google.com/view/nirvanacoppola/}

\subjclass[2020]{Primary: 11G10, Secondary: 11F11, 11F80, 14G35, 14Q15}

\begin{document}

\begin{abstract}
It is well known that the rational torsion of an abelian variety defined over a number field injects into the reduction modulo any sufficiently large prime, so the order of the torsion group divides the greatest common divisor of the sizes of points on the reduction at each prime. Drawing inspiration from Katz's Inventiones paper (1981), we investigate the converse to this for abelian varieties of \(\rm GL_2\)-type and exhibit a conjectural list of possible torsion orders for modular abelian varieties over \(\mathbb Q\) of dimension up to \(5\).
\end{abstract}

\maketitle

\section{Introduction}\label{sec:intro}
For any elliptic curve \(E\) over \(\Q\), the structure of the torsion subgroup \(\Tors(E)(\Q)\) is well understood, thanks to the celebrated theorem of Mazur \cite[Theorem 8]{Mazur1977ModularIdeal}, where he determined which finite abelian groups \(\Tors(E)(\Q)\) is isomorphic to and showed that \(|\Tors(E)(\Q)| \leq 16\).
Mazur's work was then generalised to other number fields: Merel \cite{Merel1996BornesNombres} showed that for any number field \(K\) and any elliptic curve \(E/K\) the order of the torsion subgroup \(\Tors(E)(K)\) is bounded by a constant that only depends on \(K\). In unpublished work, Oesterl\'e \cite{Oesterle1994TorsionNombres} showed that this bound can be taken to be \((3^{[K:\Q]+1})^2\) (for a proof, see \cite[Appendix A]{Derickx2016TorsionPublications}) .

For higher dimensional abelian varieties over general number fields, although the growth of the torsion group can be bounded by a constant depending on the variety and the degree of the field of definition (see e.g.\ \cite{Fourn2022TorsionField} and \cite{Clark2008LocalVarieties} for the case of local fields), no uniform bound is known.
Even for abelian surfaces, a uniform bound on the size of \(\Tors(A)(\Q)\) was only provided recently, under some assumptions on the geometric endomorphism ring \(\End_{\overline{\Q}}(A)\) of \(A\). Namely, in \cite{Laga2024RationalMultiplication}, it is shown that if \(A\) is an abelian surface with \emph{potential quaternionic multiplication} (that is if \(\End_{\overline{\Q}}(A)\) is a maximal order in a division quaternion algebra over \(\Q\)) then \(|\Tors(A)(\Q)| \leq 18\). Moreover, the only prime numbers dividing the order of a rational torsion point are \(2\) and \(3\), and there is an explicit list of possible candidates that can appear as torsion subgroups. This list gives a complete classification of such groups in the special case of abelian surfaces \emph{of \(\GL_2\)-type} with potential quaternionic multiplication.
%, i.e.\ such that the endomorphism algebra \(\End^0(A)\) contains a number field \(F\) of degree, in this case, \(2\) and, in general, degree \(\dim(A)\).

In this paper, we work more generally on abelian varieties of \(\GL_2\)-type defined over a number field \(K\), without further restrictions on the dimension and on the endomorphisms.
Our strategy is to use a classical approach in the study of rational points, that is, answering to a \emph{local-global problem}.
Indeed, for any abelian variety \(A/K\), if \(\p\) is a prime of good reduction for \(A\), the group \(\Tors(A)(K)\) injects into the finite group of \(\F_\p\)-points of the reduction of \(A\) modulo \(\p\). Thus we have that \(|\Tors(A)(K)|\) divides \(N(\p):=|\tilde{A}(\F_\p)|\), where \(\tilde{A}\) is the reduction of \(A\) modulo \(\p\).
A famous question by Lang asks if the converse is also true. More precisely he asked the following.
\begin{question}\label{question}
    Let \(A\) be an abelian variety over a number field \(K\) and let \(m\ge 2\) be an integer number. Suppose that 
    \[
    N(\p) \equiv 0 \pmod m
    \]
    holds for a set of places \(\p\) of density one. Does there exist a \(K\)-isogenous abelian variety \(A'\) for which the set of \(K\)-rational torsion points satisfies
    \[
    |\Tors(A')(K)| \equiv 0 \pmod m ?
    \]
\end{question}
Katz \cite{Katz1981GaloisVarieties} finds a positive answer to this question for the case of elliptic curves, and to a weaker version of it for the case of abelian surfaces. More precisely: the set of primes dividing the order of the torsion subgroup of at least one variety \(A'\) that is isogenous to \(A\), and the set of primes dividing all of the \(N(\p)\)'s are the same. However, he proves that there are counterexamples for abelian varieties of higher dimension, even for this weaker result.
After Katz, a lot of interest has been shown towards the classification of more counterexamples, see for example \cite{Cullinan2007LocalglobalVarieties, Cullinan2010PointsZarhin} and \cite{Cullinan2022DivisibilityFields}.
This problem is linked to the so-called \emph{local-global divisibility problem} in algebraic groups (see \cite[\S 8]{Dvornicich2022Local-globalGroups}), to which Gillibert and Ranieri give an answer in \cite{Gillibert2017OnVarieties, Gillibert2020OnVarieties}, in the case of \(\GL_2\)-type abelian varieties.

In this work, we prove the following theorem, which gives a positive answer to Katz's weaker question for \(\GL_2\)-type abelian varieties, in all dimensions.
\begin{theorem*}[{See Corollary \ref{cor:main}}] \label{thm:main2}
    Let \(A\) be an abelian variety of \(\GL_2\)-type over a number field \(K\). Let \(\Sigma\) be the set of primes \(\p\) of \(K\) of good reduction for \(A\) such that the absolute ramification index of \(\p\) is smaller than \(p-1\), where \(p\) is the rational prime under \(\p\). Then, the two integers
    \[
    \sup_{A'\text{ K-isogenous to }A} |\Tors(A')(K)|\qquad \text{ and } \qquad \gcd_{\p \in \Sigma} N(\p)
    \]
    are divisible by the same primes of good reduction for \(A\). More precisely, for any rational prime \(\ell\) of good reduction for \(A\), one has
    \[
    v_\ell\left(\sup_{A'\text{ K-isogenous to }A} |\Tors(A')(K)|\right)\le v_\ell\left(\gcd_{\p \in \Sigma} N(\p)\right).
    \]
    Moreover, if \(\ell\) is any such prime and it is totally inert in \(F\), then
    \[
    v_\ell\left(\sup_{A'\text{ K-isogenous to }A} |\Tors(A')(K)|\right)=v_\ell\left(\gcd_{\p \in \Sigma} N(\p)\right).
    \]
\end{theorem*}
This theorem will follow from Theorem \ref{thm:main_new}, where we reformulate Katz's problem in terms of Galois representations.

Finally, we show how one can apply our main results to study the torsion of \emph{modular abelian varieties} defined over the rationals, which, by Serre's conjectures, correspond to weight \(2\) newforms. 
We implement our results in the computational algebra system \texttt{Magma}, in order to produce tentative lists of possible torsion orders for modular abelian abelian varieties of dimension up to \(5\), a part of which is presented in Conjectures \ref{conjecture1} and \ref{conjecture2}.
The code and the outputs can be found at \cite{Alessandri2026MagmaVarieties}. Although it is not currently known if the lists are exhaustive, we remark that all the torsion orders appearing in the lists occur for at least one abelian variety, which can be exhibited in terms of the corresponding newform. This data enhances the data that is currently available in \texttt{LMFDB} \cite{LMFDBDatabase} and in \cite{Nicholls2018DescentCurves}. 
For example, there exists at least one abelian surface in the isogeny class associated to the level \(39\) newform with \texttt{LMFDB} label \texttt{\href{https://www.lmfdb.org/ModularForm/GL2/Q/holomorphic/39/2/a/b/}{39.2.a.b}} which
has torsion order equal to \(28\). This abelian variety is not listed in \texttt{LMFDB} and there are no abelian surfaces with torsion order equal to \(28\) in the database either.


\section{Setting}\label{sec:setting}
In this section, we recall the definition of \(\GL_2\)-type abelian variety and properties of the Galois representations attached to it, in order to reformulate Katz's problem in terms of Galois-stable lattices.

Let \(A\) be an abelian variety of dimension \(g\) defined over a number field \(K\), with fixed algebraic closure \(\bar{K}\). 
\begin{definition}\label{def:GL2}
    We say that \(A\) is of \(\GL_2\)-type if there exists a number field \(F\) of degree \(g\) over \(\Q\), with an embedding:
    \[
    F \hookrightarrow \End_K(A) \otimes_\Z \Q.
    \]
\end{definition}


\begin{remark}
    In \cite{Ribet1976GaloisMultiplications}, the author deals with abelian varieties \(A/K\) of \(\GL_2\)-type having full endomorphism algebra defined over \(K\), claiming that \emph{all of the ``serious'' results have to do} with this case. We will not make such an assumption, since our main result holds for any abelian variety of \(\GL_2\)-type.
\end{remark}

Let \(\ell\) be a rational prime. One can define the \(\ell\)-adic Galois representation attached to \(A\) as
\[
\rho_\ell : G_K = \Gal(\overline{K}/K) \rightarrow \Aut(T_\ell(A)) \subseteq \Aut(V_\ell(A)),
\]
where:
\begin{itemize}
    \item \(T_\ell(A)\) is the \(\ell\)-adic Tate module of \(A\) (see e.g. \cite{Serre1968GoodVarieties} for the definition);
    \item \(V_\ell(A)= T_\ell(A) \otimes_{\Z_\ell} \Q_\ell\).
\end{itemize}

Since \(A\) is of \(\GL_2\)-type, one has a \(F\)-rational system of compatible representations \((\rho_\lambda)_{\lambda \mid \ell}\) (see \cite{Ribet1976GaloisMultiplications, Lombardo2016ExplicitGL2-varieties, Azon2024ConductorCurves} for the construction), such that $$\rho_\lambda : G_K \rightarrow \GL_2(F_\lambda)$$ is a 2-dimensional representation. The sketch of the construction is as follows: we have \(F_\ell = F \otimes_\Q \Q_\ell = \prod_{\lambda \mid \ell} F_\lambda\), and \(V_\lambda (A) = V_\ell (A) \otimes_{F_\ell} F_\lambda\). The action of \(G_K\) on \(V_\lambda (A)\) is the representation \(\rho_\lambda\). We let \(\f_\lambda =[F_\lambda : \Q_\ell]\), and \(\mathbb{F}_\lambda\) be the residue field \(\mathcal{O}_{F_\lambda} / \lambda \mathcal{O}_{F_\lambda}\), which has size \(\ell^{\f_\lambda}\).


Moreover, we will need to use the lattice \(T_\lambda (A)= T_\ell (A) \otimes_{\mathcal{O}_{F_\ell}} \mathcal{O}_{F_\lambda}\), where \( \mathcal{O}_{F_\ell} = \mathcal{O}_F \otimes_\Z \Z_\ell\).

\begin{notation}
    Let \(\p\) be a finite place of \(K\), such that \(A\) has good reduction at \(\p\). We denote by \(\Frob_\p\) a Frobenius element in \(G_K\) corresponding to \(\p\); since \(A\) has good reduction at \(\p\), the representations \(\rho_\ell\) and \(\rho_\lambda\) are unramified at \(\p\), thus \(\rho_\ell(\Frob_\p)\) and \(\rho_\lambda(\Frob_\p)\) are well defined.
    We denote by \(P_\p^\ell(t)=\det(1-t\rho_\ell(\Frob_\p))\) the usual characteristic polynomial of Frobenius, and by \(P_\p^\lambda(t)=\det(1-t\rho_\lambda (\Frob_\p))\) the characteristic polynomial of \(\rho_\lambda(\Frob_\p)\).
\end{notation}

\textit{A priori}, \(P_\p^\lambda(t)\) has coefficients in \(F_\lambda\). However, by \cite[Theorem 2.1.2]{Ribet1976GaloisMultiplications}, the system \(\{\rho_\lambda\}_{\lambda}\) is strictly compatible and \(F\)-rational with exceptional set equal to the set of places of \(K\) at which \(A\) has bad reduction, so, with possibly these finitely many exceptions, in fact \(P_\p^\lambda (t) \in F[t]\).

Throughout the paper, for \(x \in F_\lambda\) (or \(F\)), we write
\[
x \equiv 0 \pmod{\lambda^n}
\]
whenever the \(\lambda\)-adic valuation of \(x\) is at least \(n\).

\subsection{Katz and Lang's problem}\label{sec:Katz}

For a given abelian variety \(A/K\), and \(\p\) a finite place of \(K\) of good reduction for \(A\), we denote by \(\F_\p\) the residue field of \(K\) at \(\p\), by \(\tilde{A}(\F_\p)\) the \(\F_\p\)-points on the reduction of \(A\) modulo \(\p\) and we let \(N(\p) = |\tilde{A}(\F_\p)|\).

In \cite{Katz1981GaloisVarieties}, the following question is proposed: given an abelian variety \(A/K\), and \(m\ge 2\) an integer number, if 
\[
N(\p) \equiv 0 \pmod m
\]
holds for a set of places \(\p\) of density one, does there exist an isogenous abelian variety \(A'\) for which the set of \(K\)-rational torsion points satisfies
\[
|\Tors(A')(K)| \equiv 0 \pmod m ?
\]

In op.\ cit.\ this problem is reformulated in terms of power-of-a-prime isogenies and the Galois representation \(\rho_\ell\). In particular, the question becomes the following.
Let \(A/K\) be an abelian variety and \(\ell\) a rational prime such that there exists \(n \ge 1\) satisfying
\[
\det(1-\rho_\ell(\gamma)) \equiv 0 \pmod {\ell^n }
\]
for all \(\gamma \in G_K\). Note that, by \v{C}ebotarev Density Theorem \cite[Theorem 13.4]{Neukirch1999AlgebraicTheory}, this condition is equivalent to requiring the congruence for \(\gamma\) of the form \(\Frob_\p\) for a density one set of places, namely those of good reduction not dividing \(\ell\).
Does there exist a pair of \(G_K\)-stable lattices \(\mathcal{L'} \subseteq \mathcal{L} \subseteq V_\ell(A)\) such that the quotient has order \(\ell^n\) and such that \(G_K\) acts trivially on it?

As mentioned, Katz finds a positive answer to this question for the case of elliptic curves \cite[Theorem 2, 2 bis]{Katz1981GaloisVarieties}, and to a weaker version of it for the case of abelian surfaces \cite[Theorem 4]{Katz1981GaloisVarieties}. We adapt Katz's approach for the proof of his Theorem 2 to the case of \(\GL_2\)-type abelian varieties, proving a refined version of his Theorem 4.

\section{Main results}\label{sec:main}

We will use the following result of Katz and Lang \cite[Theorem 1 (bis)]{Katz1981GaloisVarieties} on \(2\)-dimensional representations with coefficients in a local field.

\begin{theorem}\label{thm:Katz}
    Let \(L\) be a field which is complete under a discrete valuation, and whose residue field is finite. Let \(\mathcal{O}_L\) denote the ring of integers in \(L\), and let \(\pi\) denote a uniformising parameter. Let \(V\) be a 2-dimensional \(L\)-vector space, \(G \subseteq \Aut_L(V)\) a compact subgroup, and \(n \ge 1\) an integer. Suppose that we have the congruence
    \[
    \det(1-g) \equiv 0 \mod \pi^n \quad \forall g \in G.
    \]

    Then there exists a \(L\)-basis \(\{e_1,e_2\}\) of \(V\) and integers \(a,b \ge 0\) with \(a+b=n\) such that the elements of \(G\) expressed as matrices with respect to this basis lie in the subgroup
    \[
    \left\lbrace \mat{1 + \pi^a \mathcal{O}_L}{\pi^a \mathcal{O}_L}{\pi^b \mathcal{O}_L}{1 + \pi^b \mathcal{O}_L} \right\rbrace
    \]
    of \(\GL_2(\mathcal{O}_L)\).
\end{theorem}

We remark that, in \cite{Katz1981GaloisVarieties}, a detailed proof of the case where \(L=\Q_p\) is provided. However, the proof adapts immediately to the general case. We are now ready to state our main result.

\begin{theorem}\label{thm:main_new}
    Let \(A\) be an abelian variety over a number field \(K\) of \(\GL_2\)-type. Let \(\ell\) be a rational prime and let \(\lambda\) be a prime of \(F\) over \(\ell\). Suppose that, for an integer \(n \ge 1\), the congruence
    \[
    P_\p^\lambda(1) \equiv 0 \pmod {\lambda^n}
    \]
    holds for a density one set of finite places \(\p\) of \(K\). Then there exists an abelian variety \(A'\), \(K\)-isogenous to \(A\) through a \(\ell\)-power isogeny, such that 
    \[
    |\Tors(A')(K)| \equiv 0 \pmod {\ell^{\f_{\lambda}n}},
    \]
    where \(\f_\lambda\) is the inertia degree of \(\lambda \mid \ell\).
\end{theorem}
\begin{proof}
    With the notation of Theorem \ref{thm:Katz}, let \(L=F_\lambda\), so that \(\pi=\lambda\), \(V=V_\lambda (A)\) and \(G=\rho_\lambda(G_K)\). The condition \(\det(1-g)\equiv 0 \pmod {\lambda^n}\) for all \(g \in G\) is equivalent to \(P_\p^{\lambda}(1) \equiv 0 \pmod{\lambda^n}\) for a density one set of \(\p\), by \v Cebotarev Density Theorem. Then, by Theorem \ref{thm:Katz}, there exist two sublattices \(\mathcal{L'} \subseteq \mathcal{L}\) of \(V\), such that \(G_K\) acts trivially on the quotient. Such quotient has order \(\ell^{\f_\lambda n}\), indeed in the basis \(\{e_1,e_2\}\) given by the theorem, \(\mathcal{L} \cong \mathcal{O}_{F_\lambda} e_1 + \mathcal{O}_{F_\lambda} e_2\); \(\mathcal{L}' \cong \mathcal{O}_{F_\lambda} \lambda^a e_1 + \mathcal{O}_{F_\lambda} \lambda^b e_2\), with \(a+b=n\), and the quotient is
    \[
    \mathcal{L}/\mathcal{L}' \cong \dfrac{\mathcal{O}_{F_\lambda} e_1 + \mathcal{O}_{F_\lambda} e_2}{\mathcal{O}_{F_\lambda} \lambda^a e_1 + \mathcal{O}_{F_\lambda} \lambda^b e_2} \cong \dfrac{{\mathcal{O}_{F_\lambda}}}{\mathcal{O}_{F_\lambda} \lambda^a} \oplus \dfrac{\mathcal{O}_{F_\lambda}}{\mathcal{O}_{F_\lambda} \lambda^b},
    \]
    which has size \(|\mathbb{F}_\lambda|^{a+b} = \ell^{\f_\lambda n}\).

    Now, the lattice \(\mathcal{L}'\) is \(T_\lambda (A')\) for some abelian variety \(A'\) with \(\lambda\)-adic Galois representation \(\rho_\lambda\). In particular, \(A'\) is \(K\)-isogenous to \(A\) with an isogeny of degree a power of \(\ell\). Therefore, \(\Tors (A')(K)\) contains a subgroup of order \(\ell^{\f_\lambda n}\).
\end{proof}

We remark that our assumption on the congruence of \(P_\p^\lambda(1)\) modulo \(\lambda^{n}\) implies the hypothesis in Katz's problem on the number of points of the reduction of \(A\) at \(\p\). 
In fact, as already mentioned, the condition \(N(\p) \equiv 0 \pmod{\ell^n}\) is equivalent to \(P_\p^{\ell}(1)= \det(1-\rho_\ell(\Frob_\p)) \equiv 0 \pmod {\ell^n }\), since \(N(\p)=P_\p^\ell(1)\). By \cite[\S 11.9-11.10]{Shimura1967AlgebraicGroups}, the characteristic polynomial of \(\rho_\ell(\Frob_\p)\) is equal to the norm from \(F\) to \(\Q\) of the characteristic polynomial of \(\rho_{\lambda}(\Frob_\p)\), for any \(\lambda\) over \(\ell\). Then, it is a simple exercise to show that
\begin{equation} \label{eq:valuations}
    v_{\ell}(P_\p^{\ell}(1)) = v_{\ell}(\mathcal{N}_{F/\Q}(P_\p^{\lambda}(1))) = \sum_{\lambda' \mid \ell} v_{\lambda'}(P_\p^\lambda(1)) \e_{\lambda'}\f_{\lambda'},
\end{equation}
where the sum varies over all primes \(\lambda'\) lying over \(\ell\) and \(\e_{\lambda'}\) and \(\f_{\lambda'}\) are the ramification and inertia degree of \(\lambda'\), respectively. 


If \(P_\p^\lambda(1) \equiv 0 \pmod {\lambda^n}\) then \(v_{\lambda}(P_\p^\lambda(1)) \geq n\) and so \(v_{\ell}(P_\p^{\ell}(1))\) is at least \(n\e_\lambda \f_{\lambda}\). In particular, \(P_\p^{\ell}(1) \equiv 0 \pmod{\ell^n}\).

\begin{corollary}\label{cor:main}
    Let \(A\) be an abelian variety of \(\GL_2\)-type over a number field \(K\). Let \(\Sigma\) be the set of primes \(\p\) of \(K\) of good reduction for \(A\) such that the absolute ramification index of \(\p\) is smaller than \(p-1\), where \(p\) is the rational prime under \(\p\). Then, the two integers
    \[
    \sup_{A'\text{ K-isogenous to }A} |\Tors(A')(K)|\qquad \text{ and } \qquad \gcd_{\p \in \Sigma} N(\p)
    \]
    are divisible by the same primes \(\ell\) of good reduction for \(A\) (i.e. such that, for all places of \(K\) above \(\ell\), \(A\) has good reduction modulo such places). More precisely, for any rational prime \(\ell\) of good reduction for \(A\), one has
    \[
    v_\ell\left(\sup_{A'\text{ K-isogenous to }A} |\Tors(A')(K)|\right)\le v_\ell\left(\gcd_{\p \in \Sigma} N(\p)\right).
    \]
    Moreover, if \(\ell\) is any such prime and it is totally inert in \(F\), then
    \[
    v_\ell\left(\sup_{A'\text{ K-isogenous to }A} |\Tors(A')(K)|\right)=v_\ell\left(\gcd_{\p \in \Sigma} N(\p)\right).
    \]
\end{corollary}

\begin{proof}
First, let \(A'\) be an abelian variety that is \(K\)-isogenous to \(A\). By injectivity of the reduction map, which holds for all primes in \(\Sigma\) (see \cite[Appendix]{Katz1981GaloisVarieties}), \(\Tors(A')(K)\) injects into the \(\F_\p\)-points of \(A'\) for all \(\p \in \Sigma\), thus \(|\Tors(A')(K)|\) divides \( |A'(\F_\p)|=N(\p)\). Taking the sup over all \(A'\) that are \(K\)-isogenous to \(A\) and the gcd of all \(N(\p)\) as \(\p\) varies in \(\Sigma\), we obtain that
\[
    \sup_{A' K\text{-isogenous to} A} |\Tors(A')(K)| \text{ divides } \gcd_{\p \in \Sigma} N(\p).
    \]
 
Now, let \(\ell\) be a prime of good reduction for \(A\), dividing all \(N(\p)\) for \(\p \in \Sigma\), except those lying above \(\ell\). As shown in Equation \eqref{eq:valuations}, for every such \(\p\),
\begin{equation} 
    v_{\ell}(P_\p^{\ell}(1)) = \sum_{\lambda' \mid \ell} v_{\lambda'}(P_\p^\lambda(1)) \e_{\lambda'}\f_{\lambda'}=\sum_{\lambda' \mid \ell} v_{\lambda'}(P_\p^{\lambda'}(1)) \e_{\lambda'}\f_{\lambda'}.
\end{equation}

The last equality follows from the compatibility of the representations \(\{\rho_\lambda\}_\lambda\), which holds in our setting by \cite[Theorem 2.1.2]{Ribet1976GaloisMultiplications}. Since this quantity is positive, there exists \(\lambda_0\) lying over \(\ell\) with \(v_{\lambda_0}(P_\p^{\lambda_0}(1)) =n> 0\). By Theorem \ref{thm:main_new}, there exists an abelian variety \(A'\) that is \(K\)-isogenous to \(A\) and such that \(v_\ell(|\Tors(A')(K)|) \ge n \e_{\lambda_0} \f_{\lambda_0} > 0\), concluding the proof of the first statement.

Now, suppose that \(\lambda_0\) is the unique prime lying over \(\ell\), and that \(\ell\) is totally inert, i.e. \(\e_{\lambda_0}=1\). It is enough to show that there exists an abelian variety \(A'\) that is \(K\) isogenous to \(A\) and such that \(v_\ell(\gcd N(\p)) \le  v_\ell(|\Tors(A')(K)|)\). We take \(A'\) as in the previous step of the proof. Then
\[
v_\ell(N(\p))=\sum_{\lambda' \mid \ell} v_{\lambda'}(P_\p^{\lambda'}(1)) \e_{\lambda'}\f_{\lambda'} = v_{\lambda_0}(P_\p^{\lambda_0}(1))\e_{\lambda_0}\f_{\lambda_0}=n \f_{\lambda_0}.
\]

So, \(v_\ell(\gcd N(\p)) \le n \f_{\lambda_0} \le v_\ell(|\Tors(A')(K)|)\).
\end{proof}

\begin{remark}
    Notice that the set \(\Sigma\) is infinite, since all but finitely many primes are of good reduction and unramified (so in particular \(\e_p=1\)). However, in order to compute \(\gcd_{\p \in \Sigma} N(\p)\), only finitely many of these primes will be sufficient. See Lemma \ref{lem:sturm} for more details when \(K=\Q\).
\end{remark}

\section{Modular abelian varieties defined over \texorpdfstring{$\Q$}{Q}}\label{sec:app}

We now specialise the result of Theorem \ref{thm:main_new} and Corollary \ref{cor:main} to simple \(\GL_2\)-type abelian varieties defined over the rational numbers. Let \(A\) be an abelian variety of dimension \(g\) defined over \(\Q\). Let \(\End^0_\Q (A) = \End_\Q(A) \otimes_\Z \Q\) denote the algebra of endomorphisms of \(A\) that are defined over \(\Q\), and let \(\End_{\overline{\Q}}^0(A) = \End_{\overline{\Q}} (A \times_\Q \overline{\Q}) \otimes_\Z \Q\) denote the full endomorphism algebra.

By \cite[p.\ 193]{Mumford1970AbelianVarieties}, if \(A\) is a geometrically simple abelian variety, both \(\End^0_\Q (A)\) and \(\End^0_{\overline \Q} (A)\) are Albert algebras. By \cite[Remark p.\ 182]{Mumford1970AbelianVarieties}, the dimensions of \(\End^0_\Q (A)\) and of \(\End_{\overline{\Q}}^0(A)\) divide \(2g\) and by \cite[\S 2]{Ribet2004AbelianForms}, if \(\End_\Q^0(A)\) contains a number field, then its degree divides \(g\).
In particular, by Albert's classification \cite[p.\ 201]{Mumford1970AbelianVarieties}, both endomorphism algebras are of one of the following types:
\begin{itemize}
    \item \textbf{Type I}: a totally real number field of degree dividing \(g\);
    \item \textbf{Type II or III}: a quaternion algebra (either totally definite or totally indefinite) over a totally real number field of degree dividing \(g/2\);
    \item \textbf{Type IV}: a division algebra of dimension \(d^2\) over a CM field of degree dividing \(2g/d^2\).
\end{itemize}

We now investigate the possible geometric behaviours of simple abelian varieties \(A\) of \(\GL_2\)-type over \(\Q\). As in \S \ref{sec:setting}, we write \(F \hookrightarrow \End_\Q^0(A)\), where \(F\) is a degree \(g\) number field. By \cite[Theorem 2.1]{Ribet2004AbelianForms}, \(F=\End_\Q^0(A)\).

\begin{proposition}
    Let \(A\) be a simple abelian variety of dimension \(g \ge 2\), defined over \(\Q\), of \(\GL_2\)-type. Then one of the following occurs:
    \begin{itemize}
        \item \(A\) is not geometrically simple;
        \item \(A\) is geometrically simple and has potential quaternionic multiplication (or PQM for short);
        \item \(A\) is geometrically simple, with \(\End^0_{\overline{\Q}}(A) = \End^0_\Q(A)\) equal to a totally real number field of degree \(g\).
    \end{itemize}
\end{proposition}


\begin{proof}
Let us assume that \(A\) is geometrically simple, so that \(\End^0_{\overline{\Q}} (A)\) is an Albert algebra. In particular, it is either a division algebra or a field. Let \(Z\) be its center. By \cite[Proposition 5.1, Proposition 5.2, Theorem 5.3 and Proposition 3.6]{Ribet2004AbelianForms}, \(Z\) is totally real, contained in \(\End^0_{\Q} (A)\), and \(\End^0_{\overline{\Q}} (A)\) is either a quaternion algebra over \(Z\) (PQM case), or it is equal to \(Z\), in which case \(Z= \End^0_{\overline{\Q}} (A)= \End^0_{\Q} (A)\) as claimed.
\end{proof}

In the non-geometrically simple case, the study of \(\Tors(A)\) may be reduced to understanding the torsion of smaller dimensional abelian varieties, albeit defined over larger number fields. In the PQM case, for \(A\) of dimension \(2\), \cite{Laga2024RationalMultiplication} classifies all possible torsion subgroups. These first two cases are somewhat special: the generic case is that of \(\End^0_{\overline{\Q}}(A) = \End^0_\Q(A)\). Nonetheless, for what follows, we need not distinguish among these cases.


By \cite[Theorem 4.4]{Ribet2004AbelianForms} and Serre's modularity conjectures \cite{Serre1987SurmathrmGaloverlinemathbfQ/mathbfQ} (proved by Khare--Winterberger, \cite{Khare2009Serresi,Khare2009Serresii}), a simple abelian variety \(A/\Q\) of \(\GL_2\)-type is modular, hence it is isogenous to the abelian variety \(A_f\) associated to some newform \(f\), with weight \(2\) and Nebentypus \(\chi\), and coefficient field equal to \(F=\End^0_\Q(A_f)\). The definition of \(A_f\) is  in \cite[Definition 6.6.3]{Diamond2005AForms}). Such an abelian variety has \(\lambda\)-adic Galois representation isomorphic to that arising from the newform \(f\), which we denote by \(\rho_{f,\lambda}\). For the construction see \cite[\S 9.5]{Diamond2005AForms}. By \cite[Theorem 9.6.5]{Diamond2005AForms}, the representations \((\rho_{f,\lambda})\) form a strictly compatible system for all \(\lambda\)'s, \emph{including those of bad reduction for \(A_f\)}, with characteristic equation of the Frobenius attached to the rational prime \(p\) as follows:
\[
P_p^\lambda(t) = 1 - a_p(f)t + \chi(p) p t^2,
\]
where \(a_p(f)\) is the \(p\)-th Fourier coefficient of \(f\).


%which in terms of the associated modular form is equivalent to \(f\) not having non-trivial inner twists if \(A \sim A_f\), see \cite[Theorem 5.1 and \S 3]{Ribet1980TwistsVarieties}. 


This discussion allows us to restate Theorem \ref{thm:main_new} and Corollary \ref{cor:main} in the following form. The proofs are identical to those in \S \ref{sec:main}, except that we can drop the hypothesis on good reduction of \(A_f\) at \(\ell\), thanks to the compatibility of the \(\rho_{f,\lambda}\)'s.

\begin{theorem}\label{thm_mainQ}
    Let \(A/\Q\) be an abelian variety of \(\GL_2\)-type, such that \(A\) is isogenous to \(A_f\) for a modular form \(f\) with totally real coefficients field \(F\). Let \(\ell\) be a rational prime and \(\lambda\) be a prime of \(F\) lying over \(\ell\). Suppose that, for an integer \(n \ge 1\), the congruence
    \[
    \det(1-\rho_{f,\lambda}(\Frob_p)) \equiv 0 \pmod{\lambda^n}
    \]
    holds for a density one set of primes \(p\). Then, there exists an abelian variety \(A'\) that is \(\Q\)-isogenous to \(A\) through a \(\ell\)-power isogeny, such that
    \[
    |\Tors(A')(\Q)| \equiv 0 \pmod{\ell^{f_\lambda n}},
    \]
    where \(f_{\lambda}\) is the inertia degree of \(\lambda \mid \ell\).
\end{theorem}
\begin{corollary}\label{cor:mainQ}
    Let \(A\sim A_f\) be as above. Let \(\Sigma\) be the set of primes of good reduction for \(A\). Then, the two integers
    \[
    \sup_{A'\, \Q\text{-isogenous to }A_f} |\Tors(A')(\Q)|\quad \text{and} \quad \gcd_{p \in \Sigma} N(p)
    \]
    are divisible by the same primes. More precisely, for any rational prime \(\ell\), one has
    \[
    v_\ell\left(\sup_{A'\, \Q\text{-isogenous to }A} |\Tors(A')(\Q)|\right)\le v_\ell\left(\gcd_{p \in \Sigma} N(p)\right).
    \]
    Moreover, if \(\ell\) is any such prime and it is inert in \(F\), then
    \[
    v_\ell\left(\sup_{A'\, \Q\text{-isogenous to }A_f} |\Tors(A')(\Q)|\right)=v_\ell\left(\gcd_{p \in \Sigma} N(p)\right).
    \]
\end{corollary}

\subsection{Implementation and conjectural results}
In \cite{Alessandri2026MagmaVarieties}, we implement in \texttt{Magma} the results of Theorem \ref{thm_mainQ} and Corollary \ref{cor:mainQ} to produce a tentative list of torsion orders for simple modular abelian varieties of dimension \(g\) over \(\Q\). To this end, we access all\footnote{For information on the completeness of the data, see \url{https://www.lmfdb.org/ModularForm/GL2/Q/holomorphic/Completeness}} weight \(2\) newforms with level up to \(10000\) in \texttt{LMFDB} and list, for each, the following data:
\begin{itemize}
    \item the level \(N\);
    \item the Fourier coefficients \(a_p\) for \(p\) prime up to the Sturm bound (see \cite[Corollary 9.20]{Stein2007});
    \item the Nebentypus \(\chi\);
\end{itemize}
We then use these data to compute, for the non-CM forms (in particular, the associated abelian varieties are non-CM), \(P_p(1)=P_p^\lambda(1) = 1 - a_p + p\chi(p)\) for \(p\) up to the Sturm bound and such that \(\e_p < p-1\), so that torsion injects into the reduction modulo \(p\). Recall \(P_p^\lambda(t)\) does not depend on \(\lambda\), so we will drop it from the notation. We consider the factorisation in \(\Z\) of the norm from \(F\) to \(\Q\) of \(P_p(1)\) and, for each rational prime \(\ell\) occurring in the factorisation, we compute the sup, for \(\lambda\) above \(\ell\), of \(\ell^{\f_\lambda n}\), where \(\f_\lambda\) is the inertia degree of \(\lambda\) and \(n\) is the \(\lambda\)-adic valuation of \(P_p(1)\). Multiplying these terms together, we obtain the \emph{predicted torsion order},
%The idea is to loop over all newforms of weight \(2\) and level up to a certain bound, consider the newforms having dimension (i.e.\ degree of the coefficient field over \(\Q\)) equal to \texttt{deg }$=g$, and compute, for each, the greatest common divisor of \(P_p^\lambda(1)=\det(1-\rho_{f,\lambda}(\Frob_p))\), as \(p\) varies among the primes of good reduction for \(A\) between \(2\) and a certain upper bound \texttt{st}, such that \(\e_p < p-1\). We set \texttt{st} equal to the Sturm bound, 
%Ideally, we include enough primes so that the g.c.d.\ of the finitely many \(P_p^\lambda(1)\)'s we consider is equal to the actual g.c.d.\ computed using all primes. We suspect that \texttt{st} equal to the Sturm bound would suffice, and we have computational evidence to support this idea, but could not find a reference for that.
i.e.\ the product
\[
\prod_{\ell \text{ prime}} \ell^{\sup_{\lambda \mid \ell} (\f_\lambda n)},
\]
which is the value for \(\displaystyle\sup_{A'\, \Q\text{-isogenous to }A_f} |\Tors(A')(\Q)|\) predicted by Theorem \ref{thm_mainQ}. Moreover, the program checks whether the integer \(\gcd_p N(p)\) (where the g.c.d.\ is over all the primes \(p\) in the aforementioned loop) is equal to the \emph{predicted torsion order}, in which case we know that there cannot be an abelian variety in the same isogeny class with torsion group of a larger size. %, since there is an injection of \(\Tors(A')(\Q)\) into all the \(A'(\F_p)\).
The output of the code is 
\begin{itemize}
    \item the list of all possible torsion orders found in this way, 
    \item the list of the torsion orders which are equal to the corresponding integer \(\gcd_p N(p)\),
    \item the list of primes which occur in the factorisation of at least one of the possible torsion orders found. 
\end{itemize}
The \emph{predicted torsion order} is not always equal to the greatest common divisor of the \(N(p)\)'s, hence Theorem \ref{thm_mainQ} and Corollary \ref{cor:mainQ} cannot be improved to be exactly equivalent to Katz's result for elliptic curves. Nonetheless, this is often the case and, using all the data in \texttt{LMFDB}, the two lists computed by the code are exactly the same for abelian varieties of dimension \(2\), \(3\) and \(5\).
In dimension \(4\), the newform with \texttt{LMFDB} label \texttt{\href{https://www.lmfdb.org/ModularForm/GL2/Q/holomorphic/4830/2/a/ce/}{4830.2.a.ce}} has an associated abelian variety (up to isogeny) with \emph{predicted torsion order} equal to \(88\), and \(\gcd_p N(p) = 176\), but \(88\) does not appear in the list of the \(\gcd_p N(p)\)'s, while \(176\) does appear in the list of predicted torsion orders. It is unclear whether, allowing the code to search for higher level newforms, this will remain an exception (and, if so, the only one) or not.

In the following lemma, we prove that using only the Fourier coefficients \(a_p\) with \(p\) up to the Sturm bound ensures the correctness of the predicted torsion orders.

\begin{lemma}\label{lem:sturm}
    Let \(f\) be a newform of weight \(2\), level \(N\) and Nebentypus \(\chi\). Let \texttt{st} be the Sturm bound for the space of such modular forms. Then, the ideal generated by \(a_p(f) - (1 + \chi(p)p)\) for every \(p\) is the same as the ideal generated by the same elements, for \(p\) up to \texttt{st}.
\end{lemma}

\begin{proof}
    Let \(g\) be an Eisenstein series having \(a_p(g)=1+\chi(p)p\) for every prime \(p\) not dividing the level. This exists, since it can be written as a linear combination of the series in the basis found in \cite[\S 4.6]{Diamond2005AForms} (taking \(\varphi=\chi,\psi=1\) and renormalising the coefficients). Let \(h=f-g\). Then for every prime \(p\) not dividing \(N\), 
    \[
    a_p(h) = a_p(f) - (1 + \chi(p)p).
    \]
    The claim is equivalent to saying that all the \(a_p(h)\) are in the ideal generated by the first \(a_p(h)\) with \(p\) up to \texttt{st}. This follows immediately from the Sturm theorem \cite[Corollary 9.20]{Stein2007} applied to the modular form \(h\).
\end{proof}

We conclude by stating two conjectures, which are suggested by the outcome of the code, on the possible torsion orders of non-CM simple abelian varieties of \(\GL_2\)-type and dimension up to \(5\), and on the primes dividing such torsion orders. All the numbers in the table are divisors of the predicted torsion orders computed by the code, and the list of newforms with the respective predicted torsion orders can be found at \cite{Alessandri2026MagmaVarieties}. We do not currently know whether the lists in Conjectures \ref{conjecture1} and \ref{conjecture2} are exhaustive. Code and data on how to associate in practice a classical modular form of weight \(2\), dimension \(2\), and trivial Nebentypus to all genus \(2\) curves of \(\GL_2\)-type whose Jacobians are \(\Q\)-isogenous to the modular abelian varieties associated to \(f\) is available at \cite{CostaModularAbelianSurfaces}.



\begin{conjecture}\label{conjecture1}
    Let \(A/\Q\) be a simple non-CM \(g\)-dimensional abelian variety of \(\GL_2\)-type. The possible orders of \(\Tors(A)(\Q)\) are listed below.
    \begin{center}
        \begin{table}[h!]
            \centering
            \begin{tabular}{r|l}
                \(g\) & \(|\Tors(A)(\Q)|\) \\
                \hline
                \(2\) & \(1, 2, 3, 4, 5, 6, 7, 8, 9, 10, 11, 12, 13, 14, 15, 16, 17, 
18, 19, 20, 21, 22, 23, 24, 28,\)\\ 
&\(31, 37, 44, 56\) \\
                \(3\) & \(1, 2, 3, 4, 5, 6, 7, 8, 9, 10, 11, 12, 13, 14, 15, 16, 17, 
18, 19, 20, 22, 23, 24, 28, 29,\)\\
&\(31, 32, 36, 37, 38, 40, 44, 46, 49, 58, 80, 83, 92, 160\)\\
                \(4\) & \(1, 2, 3, 4, 5, 6, 7, 8, 9, 10, 11, 12, 13, 14, 15, 16, 17, 
18, 19, 20, 21, 22, 23, 24, 25,\)\\
&\(26, 27, 28, 29, 31, 32, 34, 35, 36, 37, 38, 40, 41, 43, 44, 48, 49, 52, 55, 56, 57, 61,\)\\
&\(
63, 64, 68, 71, 72, 73, 74, 76, 80, 82, 88, 136, 146, 152, 155, 176\)\\
                \(5\) & \(1, 2, 3, 4, 5, 6, 7, 8, 9, 10, 11, 12, 13, 14, 16, 17, 18, 
19, 20, 21, 22, 23, 24, 25, 26,\)\\
&\(27, 28, 29, 31, 32, 34, 35, 36, 38, 40, 41, 43, 44, 46, 48, 50, 52, 53, 55, 56, 63, 64,\)\\
&\(68, 72, 80, 92, 96, 106, 110, 144, 160, 192, 288, 320\)\\
            \end{tabular}
            %\caption{Torsion}
            \label{tab:torsion}
        \end{table}
    \end{center}
\end{conjecture}

\begin{conjecture}\label{conjecture2}
    Let \(A/\Q\) be a simple non-CM \(g\)-dimensional abelian variety of \(\GL_2\)-type and let \(\ell\) be a prime such that \(A\) admits a \(\ell\)-torsion point.
    \begin{itemize}
        \item If \(g=2\), then \(\ell \in \{2, 3, 5, 7, 11, 13, 17, 19, 23, 31, 37\}\).
        \item If \(g=3\), then \(\ell \in \{2, 3, 5, 7, 11, 13, 17, 19, 23, 29, 31, 37, 83\}\).
        \item If \(g=4\), then \(\ell \in \{2, 3, 5, 7, 11, 13, 17, 19, 23, 29, 31, 37, 41, 43, 61, 71, 73 \}\).
        \item If \(g=5\), then \(\ell \in \{2, 3, 5, 7, 11, 13, 17, 19, 23, 29, 31, 41, 43, 53 \}\).
    \end{itemize}
\end{conjecture}



\section*{Acknowledgments}
The authors are grateful to Martin Azon, Sander Dahmen, Sam Frengley, Julian Demeio, Pip Goodman, Davide Lombardo, Laura Paladino and Matteo Verzobio for contributing to making this collaboration possible and for their useful insights.
\section*{Funding}
The authors are members of the INdAM group GNSAGA. JA was supported by UKRI Future Leaders Fellowship \texttt{MR/V021362/1} and the Max Planck Institute for Mathematics in Bonn, that she thanks for their hospitality and financial support. NC is supported by the \emph{``Piano di Sviluppo Dipartimentale''} with the title \emph{``Enhancing the Research in the Math Dept'', CUP C93C24000160005}. NC thanks University of Bath for its hospitality during her visit in May 2025, where most of this work has been carried on. 
\bibliographystyle{alpha}
\bibliography{refs}
\end{document}
